\documentclass[
  aps,
  prc,
  reprint,
  superscriptaddress,
  nofootinbib
]{revtex4-2}

% ----------------------------------------------------------------
% Packages
% ----------------------------------------------------------------
\usepackage[T1]{fontenc}
\usepackage[utf8]{inputenc}
\usepackage{lmodern}
\usepackage{amsmath,amssymb,bm}
\usepackage{graphicx}
\usepackage{booktabs}
\usepackage{siunitx}
\usepackage{microtype}
\usepackage{xcolor}
\usepackage{hyperref}

\hypersetup{
  colorlinks=true,
  linkcolor=blue,
  citecolor=blue,
  urlcolor=blue,
  pdftitle={Nuclear-State Control as a Dynamical Degree of Freedom in Fission Systems},
  pdfauthor={Author Name}
}

% ----------------------------------------------------------------
% Convenience macros
% ----------------------------------------------------------------
\newcommand{\keff}{k_{\mathrm{eff}}}
\newcommand{\sigf}{\sigma_f}
\newcommand{\Sigf}{\Sigma_f}
\newcommand{\rhoN}{\rho_N}
\newcommand{\rhoQ}{\rho_Q}
\newcommand{\Rf}{R_f}

% ----------------------------------------------------------------
% Document
% ----------------------------------------------------------------
\begin{document}

\title{Quantum-State Engineering of Fission Pathways in Uranium Fuel}

\author{Aidar S. Nurqalin}
\affiliation{JV Inkai LLP, Turkistan Region, Republic of Kazakhstan}

\author{Élise Morveaux}
\affiliation{Orano, Philippe Coste Conversion Plant, Tricastin, France}

\author{Tomas van der Meerveld}
\affiliation{Urenco Nederland B.V., Almelo, The Netherlands}

\author{Karin E. Lindqvaller}
\affiliation{Westinghouse Electric Sweden AB, Västerås, Sweden}

\author{Lukas M. Kellenwald}
\affiliation{Paul Scherrer Institut, 5232 Villigen PSI, Switzerland}

\date{\today}

\begin{abstract}
Nuclear reactor physics conventionally treats microscopic quantum dynamics through nuclear reaction data while controlling fission primarily through macroscopic variables such as neutron population, absorption, moderation, geometry, and temperature. We investigate a complementary possibility: whether the internal nuclear state of fissile matter can be promoted to an explicitly controlled dynamical variable. In this framework, an external actuator modifies a nuclear-state distribution, the state-dependent fission response changes, and the resulting perturbation propagates to the macroscopic fission rate and neutron multiplication. We formulate the concept using a density-matrix description, derive a minimal two-state control model, and connect it to an effective fission cross section and reactor reactivity. The low-lying isomer of $^{235}$U provides a concrete candidate because its neutron-induced fission cross section is experimentally known to differ from that of the ground state. Existing literature has already suggested exploiting this difference for fission control; the present contribution instead focuses on formulating nuclear-state population as a dynamical control variable and embedding it in a multiscale chain from microscopic state preparation to macroscopic reactor kinetics. The framework yields a falsifiable experimental criterion and separates three questions: state dependence, external actuation, and system-level coupling.
\end{abstract}

\keywords{nuclear fission, nuclear quantum control, reactor kinetics, nuclear isomers, $^{235m}$U, coherent control}

\maketitle

\section{Introduction}
\label{sec:introduction}

Nuclear fission is intrinsically a quantum-mechanical process. The formation of an excited compound nucleus, the evolution of nuclear degrees of freedom toward scission, quantum tunnelling through the fission barrier, and the subsequent formation of fragments are all governed, at the microscopic level, by quantum dynamics. Modern microscopic approaches therefore describe the nucleus as a complex interacting many-body quantum system and employ nuclear density functional theory, time-dependent methods, and collective quantum dynamics to investigate the path from excitation to scission \cite{SchunckRobledo2016}.

At the reactor-engineering scale, however, this microscopic quantum dynamics is not treated as an active control degree of freedom. Reactor physics is formulated primarily in terms of neutron transport, reaction rates, isotopic concentrations, thermal-hydraulic variables, and macroscopic feedback mechanisms. The underlying nuclear dynamics enters through quantities such as reaction cross sections, decay constants, neutron multiplicities, and evaluated nuclear data. Consequently, the internal quantum state of the fissile nuclei is normally not regarded as a dynamically controllable reactor variable.

Reactor control therefore acts predominantly on the environment in which fission takes place rather than on the nuclear state from which fission proceeds. Control rods modify neutron absorption; moderators alter neutron energy distributions; soluble absorbers modify macroscopic absorption; coolant and fuel temperatures provide reactivity feedback; and changes in geometry or composition affect neutron leakage and multiplication. In each case, the principal strategy is to control the neutron population or its interaction with an effectively passive fissile medium.

This situation differs from that encountered in atomic and optical physics. Coherent control makes it possible not only to predict transition probabilities but also to manipulate quantum states by controlling the amplitude, frequency, timing, and phase of applied electromagnetic fields.

Related approaches to the representation and manipulation of controlled
quantum states have also been discussed in the context of quantum-system
design and simulation \cite{SisiniPC}.

Nuclear quantum optics has progressively extended related concepts to nuclear degrees of freedom. Early theoretical work explicitly identified preparation, control, and detection of nuclear states as a possible objective of resonant laser--nucleus interactions \cite{Buervenich2006}.

Subsequent experiments have demonstrated that this idea is physically meaningful in selected systems. Coherent X-ray-optical manipulation of collective nuclear excitations has been demonstrated in $^{57}$Fe \cite{Heeg2021}. Spectroscopic access to the exceptionally low-energy isomeric transition in $^{229}$Th has further shown that nuclear transitions can be addressed with optical-frequency techniques under favorable conditions \cite{Yamaguchi2024}. At much higher photon energies, theoretical work has investigated coherent MeV photon fields interacting with heavy nuclei in regimes where excitation competes with internal equilibration, neutron evaporation, and fission \cite{Palffy2015}.

These developments motivate a question that is distinct from both conventional reactor physics and conventional nuclear spectroscopy:

\begin{quote}
Can externally controlled preparation or manipulation of nuclear states in fissile matter modify the subsequent probability or dynamics of neutron-induced fission in a way that can be represented as a controllable physical variable?
\end{quote}

The question is not whether electromagnetic radiation can itself induce fission, nor whether fission is fundamentally quantum mechanical. Rather, the issue considered here is whether the \emph{state dependence} of the fission process can be exploited as a control mechanism.

Conventionally, a fission rate may be represented schematically as
\begin{equation}
R_f
=
N_f
\int
\sigma_f(E)\,
\Phi(E)\,
dE,
\label{eq:conventional_fission_rate_intro}
\end{equation}
where $N_f$ is the number density or number of fissile nuclei appropriate to the chosen normalization, $\Phi(E)$ is the neutron flux, and $\sigma_f(E)$ is the neutron-induced microscopic fission cross section.

Within this description, reactor control primarily modifies the neutron field, while $\sigma_f(E)$ is treated as a material response supplied by nuclear data. The conceptual extension investigated in this work is to allow the effective fission response to depend explicitly on a nuclear-state distribution,
\begin{equation}
\sigma_f
\rightarrow
\sigma_f(E,\rho_N),
\label{eq:sigma_state_intro}
\end{equation}
where $\rho_N$ denotes a nuclear density operator or, in the incoherent limit, a statistical distribution over nuclear states.

If $\rho_N$ can be influenced by an external control variable $u_Q(t)$,
\begin{equation}
u_Q(t)
\longrightarrow
\rho_N(t)
\longrightarrow
\sigma_f^{\mathrm{eff}}(t)
\longrightarrow
R_f(t),
\label{eq:central_chain_intro}
\end{equation}
then the microscopic state of the fissile material becomes, in principle, an additional degree of freedom in a controlled fission system.

This possibility should not be interpreted as a proposal for an immediately realizable ``quantum nuclear reactor.'' Significant obstacles are evident, including state-preparation efficiency, competing decay channels, rapid equilibration of highly excited nuclei, decoherence, the fraction of nuclei that can be addressed, and the energetic cost of actuation. The purpose of this work is more limited and more fundamental: to formulate the problem as a consistent dynamical control framework and to identify the minimum physical conditions required for nuclear-state manipulation to influence a macroscopic fission system.

The central hypothesis may therefore be stated as
\begin{equation}
\boxed{
\begin{minipage}{0.86\columnwidth}
\centering
The nuclear state of fissile matter, normally implicit in reaction data,\\
may in principle become an explicit control variable.
\end{minipage}
}
\end{equation}

The remainder of this paper develops this statement from conventional reactor control to a minimal microscopic model, identifies experimentally established state-dependent fission in $^{235m}$U as a concrete anchor, and derives a falsifiable path from nuclear-state actuation to reactor-level observables.

\section{From Conventional Reactor Control to Nuclear-State Control}
\label{sec:control_framework}

Conventional reactor control is based on the manipulation of macroscopic variables that influence neutron balance. At the simplest level, the evolution of a multiplying system may be represented through the effective multiplication factor $k_{\mathrm{eff}}$, or equivalently through the reactivity
\begin{equation}
\rho
=
\frac{k_{\mathrm{eff}}-1}{k_{\mathrm{eff}}}.
\end{equation}

For a fissile population $N_f$, the neutron-induced fission rate can be written schematically as
\begin{equation}
R_f(t)
=
N_f
\int
\sigma_f(E)\,
\Phi(E,t)\,
dE.
\label{eq:fission_rate_control}
\end{equation}

In this formulation, the main time-dependent quantity accessible to conventional control is the neutron field. Control rods, soluble absorbers, moderator conditions, coolant density, fuel temperature, and geometry ultimately modify neutron production, absorption, moderation, or leakage.

A simplified conventional control chain may therefore be represented as
\begin{equation}
u_C(t)
\longrightarrow
\mathcal{A}_C
\longrightarrow
\Phi(E,t)
\longrightarrow
R_f(t),
\end{equation}
where $u_C(t)$ denotes a conventional command and $\mathcal{A}_C$ the associated physical actuator.

The nuclear properties entering $\sigma_f(E)$ are not normally treated as dynamically controllable quantities. For a given isotope and thermodynamic regime, they are effectively constitutive properties encoded in evaluated reaction data.

The present work investigates the consequences of relaxing this assumption.

\subsection{Nuclear state as an explicit variable}

Let the state of the fissile nuclear ensemble be represented by a density operator
\begin{equation}
\rho_N.
\end{equation}

A density operator is preferable to a pure state because macroscopic fissile material will generally contain an ensemble of nuclei occupying multiple states and interacting with a complex environment.

The conventional fission cross section may then be regarded as an ensemble-averaged response associated with an uncontrolled or reference distribution,
\begin{equation}
\sigma_f(E)
=
\sigma_f(E;\rho_N^{(0)}).
\end{equation}

If an external interaction changes the state distribution,
\begin{equation}
\rho_N^{(0)}
\rightarrow
\rho_N(t),
\end{equation}
the more general response is
\begin{equation}
\sigma_f
=
\sigma_f(E;\rho_N).
\end{equation}

The fission rate becomes
\begin{equation}
R_f(t)
=
N_f
\int
\sigma_f(E;\rho_N(t))
\Phi(E,t)\,dE.
\label{eq:fission_rate_state}
\end{equation}

Equation~\eqref{eq:fission_rate_state} contains two conceptually distinct control channels:
\begin{equation}
\Phi(E,t)
\rightarrow
R_f(t)
\end{equation}
and
\begin{equation}
\rho_N(t)
\rightarrow
\sigma_f^{\mathrm{eff}}(t)
\rightarrow
R_f(t).
\end{equation}

Conventional control modifies the neutron environment. Nuclear-state control would instead modify the response of the fissile material itself.

\subsection{A second control channel}

We introduce a generic external control variable $u_Q(t)$ representing an interaction capable of modifying the nuclear state:
\begin{equation}
u_Q(t)
\longrightarrow
\rho_N(t)
\longrightarrow
\sigma_f^{\mathrm{eff}}(t)
\longrightarrow
R_f(t).
\label{eq:q_control_chain}
\end{equation}

Its physical implementation is deliberately left unspecified. It could correspond to a coherent electromagnetic field, an electron-mediated process, an X-ray or gamma-ray pulse sequence, or another externally tunable interaction.

The complete controlled fission system may therefore possess two parallel pathways,
\begin{align}
u_C(t) &\longrightarrow \Phi(E,t),\\
u_Q(t) &\longrightarrow \rho_N(t),
\end{align}
which both influence $R_f(t)$.

\subsection{State-dependent fission probability}

Consider a set of nuclear states $\{|i\rangle\}$, each characterized by a state-dependent fission response $\sigma_f^{(i)}(E)$. For an incoherent mixture with populations $p_i$,
\begin{equation}
\sigma_f^{\mathrm{eff}}(E)
=
\sum_i
p_i
\sigma_f^{(i)}(E).
\label{eq:incoherent_sigma}
\end{equation}

Thus, external redistribution of state populations,
\begin{equation}
p_i^{(0)}
\rightarrow
p_i(t),
\end{equation}
is sufficient to modify the effective fission response if at least two states have unequal fission cross sections.

A more general quantum description contains off-diagonal density-matrix elements $\rho_{ij}$ with $i\neq j$. Formally, the fission probability can be written as
\begin{equation}
P_f(E)
=
\mathrm{Tr}
\left[
\rho_N\,
\hat{\Pi}_f(E)
\right],
\label{eq:fission_operator}
\end{equation}
where $\hat{\Pi}_f(E)$ denotes an effective operator associated with evolution into the fission channel.

Equation~\eqref{eq:fission_operator} expresses the central idea compactly: if an external control changes $\rho_N$ and the fission operator is state sensitive, then the fission probability is externally modifiable.

\subsection{From microscopic control to macroscopic reactivity}

Let
\begin{equation}
\delta\sigma_f
=
\sigma_f^{\mathrm{controlled}}
-
\sigma_f^{(0)}
\end{equation}
be the state-induced modification of the microscopic fission response. At the macroscopic level,
\begin{equation}
\delta\Sigma_f
=
N_f\,\delta\sigma_f.
\end{equation}

A nonzero $\delta\Sigma_f$ perturbs neutron production and therefore, in principle, the multiplication factor:
\begin{equation}
\delta\rho_N
\rightarrow
\delta\sigma_f
\rightarrow
\delta\Sigma_f
\rightarrow
\delta k_{\mathrm{eff}}
\rightarrow
\delta\rho.
\end{equation}

For bookkeeping, we introduce a hypothetical effective reactivity contribution $\rho_Q$ and write, at first order,
\begin{equation}
\rho_{\mathrm{tot}}
=
\rho_{\mathrm{conv}}
+
\rho_Q.
\label{eq:rho_decomposition}
\end{equation}

Equation~\eqref{eq:rho_decomposition} is not assumed to hold as an exact additive decomposition in arbitrary nonlinear regimes; it is a perturbative definition that permits direct comparison with conventional reactivity effects.

\subsection{Necessary physical conditions}

For nuclear-state control to constitute a meaningful additional degree of freedom, several conditions must hold:
\begin{enumerate}
\item an external interaction must reproducibly modify the nuclear-state distribution,
\begin{equation}
u_Q(t)\rightarrow\delta\rho_N(t);
\end{equation}
\item the relevant states must have distinguishable fission responses,
\begin{equation}
\frac{\partial P_f}{\partial\rho_N}\neq0;
\end{equation}
\item the state modification must survive long enough to affect the relevant fission interaction;
\item a sufficiently large fraction of nuclei must be addressed;
\item the ensemble-averaged change must survive competing relaxation, statistical fluctuations, and ordinary reactor feedback.
\end{enumerate}

The proposed mechanism can therefore be summarized by
\begin{equation}
\boxed{
u_Q
\rightarrow
\delta\rho_N
\rightarrow
\delta P_f
\rightarrow
\delta\Sigma_f
\rightarrow
\delta k_{\mathrm{eff}}
}.
\label{eq:necessary_chain}
\end{equation}

Every arrow in Eq.~\eqref{eq:necessary_chain} is a separate physical problem. The following sections examine which links are already supported by existing nuclear physics and which remain open.

\section{Physical Basis for Nuclear-State Control of Fission}
\label{sec:physical_basis}

The framework requires two independent physical ingredients. First, nuclear states must be externally addressable or controllable in at least some systems. Second, fission probability must depend on the populated nuclear state. Neither statement is intrinsically speculative; what remains unestablished is whether the two can be combined in a regime that yields useful, externally controlled modification of fission.

\subsection{External manipulation of nuclear states}

The first requirement is
\begin{equation}
u_Q(t)
\longrightarrow
\delta\rho_N(t).
\end{equation}

Nuclear quantum optics provides direct evidence that nuclear excitation dynamics can be manipulated by external fields. Bürvenich, Evers, and Keitel showed theoretically that high-frequency laser facilities could enable resonant laser--nucleus interactions and explicitly identified preparation and control of nuclear states as a long-term objective \cite{Buervenich2006}.

Experimentally, coherent X-ray-optical control of collective nuclear excitations has been demonstrated in $^{57}$Fe. Heeg \textit{et al.} used phase-controlled X-ray pulses to switch nuclear dynamics between enhanced coherent excitation and stimulated emission \cite{Heeg2021}. This establishes that phase-sensitive manipulation of a nuclear excitation is experimentally possible.

The $^{229}$Th system provides a complementary result. Its unusually low-energy nuclear isomeric transition lies in the optical/VUV domain, and high-resolution laser spectroscopy of the isomeric transition has now been achieved \cite{Yamaguchi2024}. Although $^{229}$Th is not itself the principal fissile candidate considered below, it demonstrates that the internal state of a nucleus is not universally inaccessible to external optical control.

At much higher excitation energies, Pálffy \textit{et al.} investigated coherent MeV photon pulses interacting with heavy nuclei in a quasiadiabatic regime where photon absorption competes with intrinsic nuclear equilibration. Their model includes neutron evaporation and fission among the decay channels \cite{Palffy2015}. This work does not demonstrate controlled fission, but it shows that externally driven coherent excitation and fission can occur within the same dynamical framework.

\subsection{Fission probability is state dependent}

The second requirement is
\begin{equation}
\frac{\partial P_f}{\partial\rho_N}
\neq0.
\end{equation}

In compound-nucleus reaction theory, decay probabilities depend not only on excitation energy but also on quantities such as angular momentum, parity, transition-state structure, and available channels. Schematically,
\begin{equation}
\sigma_{n,f}(E)
=
\sum_{J,\pi}
\sigma^{CN}(E,J,\pi)
G_f(E^\ast,J,\pi),
\end{equation}
where $\sigma^{CN}$ denotes compound-nucleus formation and $G_f$ the probability of decay through fission.

Thus, even in a statistical description,
\begin{equation}
G_f(E^\ast,J_1,\pi_1)
\neq
G_f(E^\ast,J_2,\pi_2)
\end{equation}
in general.

The strongest anchor for the present work is more direct: the low-lying isomer of $^{235}$U has a measured neutron-induced fission cross section different from that of the ground state \cite{DEer1988}. Consequently, the link
\begin{equation}
\rho_N\rightarrow\sigma_f
\end{equation}
is experimentally established in a relevant fissile isotope.

\subsection{Fission as quantum many-body dynamics}

Modern microscopic theory treats induced fission as large-amplitude, time-dependent many-body dynamics \cite{SchunckRobledo2016}. Time-dependent nuclear density-functional calculations follow the evolution toward scission and show sensitivity to the initial configuration and dynamical trajectory. Finite-temperature TDDFT studies, for example, explicitly investigate dissipation and trajectory-dependent evolution during induced fission \cite{Li2023}.

Alternative microscopic descriptions also emphasize that only subsets of many-body states may contribute strongly to transmission across the fission barrier. A nonequilibrium Green-function treatment of $^{236}$U found that the fluctuations of fission probability can be associated with the number of many-body eigenstates that simultaneously have significant support on both sides of the barrier and lie near the system excitation energy \cite{UzawaHagino2024}.

These results motivate the search for pairs or families of configurations satisfying
\begin{equation}
P_f^{(A)}
\neq
P_f^{(B)}
\end{equation}
together with a physically accessible external transition or redistribution
\begin{equation}
|\Psi_A\rangle
\overset{u_Q}{\longleftrightarrow}
|\Psi_B\rangle.
\end{equation}

\subsection{Population control versus coherent control}

Two different physical regimes should be distinguished.

In the weaker, population-based regime, the density matrix is effectively diagonal,
\begin{equation}
\rho_N
=
\sum_i
p_i
|i\rangle\langle i|,
\end{equation}
and
\begin{equation}
P_f^{\mathrm{eff}}
=
\sum_i
p_i P_f^{(i)}.
\end{equation}

Here, genuine long-lived coherence is unnecessary. Redistribution of state populations is sufficient if the states have different fission responses.

In the stronger regime, the external interaction creates a coherent superposition
\begin{equation}
|\Psi\rangle
=
\sum_i c_i|i\rangle
\end{equation}
with nonzero off-diagonal density-matrix elements. If multiple components feed common fission channels, the transition amplitude can contain interference terms. Schematically,
\begin{equation}
A_f
=
\sum_i c_i A_{i\rightarrow f},
\end{equation}
so that
\begin{equation}
P_f
=
|A_f|^2
\end{equation}
contains cross terms depending on relative phases.

A phase-dependent fission response would represent genuine coherent control, but it is also the more demanding scenario because highly excited heavy nuclei possess dense spectra and rapid internal equilibration. The overall control concept therefore does not require coherent macroscopic superpositions: population control alone is sufficient to define an additional state-dependent fission-control channel.

\subsection{The missing connection}

Existing physics supports the statements
\begin{equation}
\boxed{\text{nuclear states can be externally manipulated}}
\end{equation}
and
\begin{equation}
\boxed{\text{fission response can depend on nuclear state}}.
\end{equation}

What remains to be established in a suitable fissile system is the controlled connection
\begin{equation}
\boxed{
\begin{minipage}{0.86\columnwidth}
\centering
external actuation $\rightarrow$ controlled nuclear-state redistribution\\
$\rightarrow$ controlled modification of neutron-induced fission
\end{minipage}
}.
\end{equation}

This missing connection defines the central microscopic problem of the present framework.

\section{Minimal Theoretical Model}
\label{sec:minimal_model}

The purpose of this section is to construct the simplest dynamical model capable of representing nuclear-state control of fission. The model is intentionally abstract and does not initially assume a specific isotope.

Consider two nuclear configurations,
\begin{equation}
|A\rangle,
\qquad
|B\rangle,
\end{equation}
with different fission decay rates,
\begin{equation}
\Gamma_f^{(A)}
\neq
\Gamma_f^{(B)}.
\end{equation}

An external actuator couples the two states,
\begin{equation}
|A\rangle
\overset{u_Q(t)}{\longleftrightarrow}
|B\rangle.
\end{equation}

\subsection{Population-based two-state model}

Let $p_A(t)$ and $p_B(t)$ denote the two populations. The external actuator produces transition rates
\begin{equation}
W_{A\rightarrow B}(t),
\qquad
W_{B\rightarrow A}(t).
\end{equation}

Both states may decay through fission and through competing channels. Define
\begin{align}
\Gamma_A &= \Gamma_f^{(A)}+\Gamma_d^{(A)},\\
\Gamma_B &= \Gamma_f^{(B)}+\Gamma_d^{(B)}.
\end{align}

The population equations are
\begin{align}
\frac{dp_A}{dt}
&=
-
W_{A\rightarrow B}p_A
+
W_{B\rightarrow A}p_B
-
\Gamma_Ap_A,
\\
\frac{dp_B}{dt}
&=
W_{A\rightarrow B}p_A
-
W_{B\rightarrow A}p_B
-
\Gamma_Bp_B.
\end{align}

The instantaneous fission rate for the nuclear ensemble is
\begin{equation}
r_f(t)
=
\Gamma_f^{(A)}p_A(t)
+
\Gamma_f^{(B)}p_B(t),
\end{equation}
and the integrated fission probability is
\begin{equation}
P_f
=
\int_0^\infty r_f(t)\,dt.
\end{equation}

For an initially populated state $A$ and no external control,
\begin{equation}
P_f^{(0)}
=
\frac{\Gamma_f^{(A)}}{\Gamma_A}.
\end{equation}

If population can be transferred to state $B$, control is possible whenever the two branching ratios differ:
\begin{equation}
\boxed{
\frac{\Gamma_f^{(A)}}{\Gamma_A}
\neq
\frac{\Gamma_f^{(B)}}{\Gamma_B}
}
\end{equation}
together with
\begin{equation}
\boxed{
W_{A\rightarrow B}\neq0.
}
\end{equation}

No long-lived quantum coherence is required.

\subsection{Microscopic control efficiency}

We define a dimensionless microscopic control efficiency,
\begin{equation}
\eta_Q
=
\frac{P_f^{(Q)}-P_f^{(0)}}{P_f^{(0)}},
\end{equation}
where $P_f^{(Q)}$ denotes the fission probability under external actuation.

Positive $\eta_Q$ corresponds to enhanced fission, negative $\eta_Q$ to suppression. The quantity provides a direct metric that can be calculated microscopically or measured experimentally.

\subsection{Competition between control and relaxation}

State transfer must compete with intrinsic relaxation and equilibration. Introduce
\begin{equation}
\tau_Q\sim\frac{1}{W_Q}
\end{equation}
and
\begin{equation}
\tau_{\mathrm{eq}}
\sim
\frac{1}{\Gamma_{\mathrm{eq}}}.
\end{equation}

A necessary condition for the prepared state to influence fission is approximately
\begin{equation}
\tau_Q
\lesssim
\tau_{\mathrm{eq}},
\end{equation}
or
\begin{equation}
W_Q
\gtrsim
\Gamma_{\mathrm{eq}}.
\end{equation}

We therefore define
\begin{equation}
\mathcal{C}_Q
=
\frac{W_Q}{\Gamma_{\mathrm{eq}}},
\end{equation}
with $\mathcal{C}_Q\ll1$, $\sim1$, and $\gg1$ representing ineffective, competing, and dominant external control, respectively.

\subsection{Coherent two-state extension}

A coherent version of the model may be represented by
\begin{equation}
H(t)
=
\frac{\hbar}{2}
\begin{pmatrix}
-\Delta &
\Omega(t)e^{-i\phi(t)}
\\
\Omega(t)e^{i\phi(t)} &
\Delta
\end{pmatrix},
\end{equation}
where $\Omega(t)$ is a Rabi coupling, $\Delta$ the detuning, and $\phi(t)$ an externally controlled phase.

The density matrix evolves according to a master equation,
\begin{equation}
\frac{d\rho}{dt}
=
-\frac{i}{\hbar}[H,\rho]
+
\mathcal{L}_{\mathrm{decay}}[\rho]
+
\mathcal{L}_{\mathrm{decoh}}[\rho].
\end{equation}

The fission rate can then contain both population and interference contributions,
\begin{equation}
r_f(t)
=
\Gamma_f^{(A)}\rho_{AA}
+
\Gamma_f^{(B)}\rho_{BB}
+
r_f^{(\mathrm{int})}.
\end{equation}

If both states feed common fission channels,
\begin{equation}
r_f^{(\mathrm{int})}
\propto
2\,\mathrm{Re}
\left[
\rho_{AB}
A_{A\rightarrow f}
A_{B\rightarrow f}^{\ast}
\right].
\end{equation}

A response satisfying
\begin{equation}
\frac{\partial P_f}{\partial\phi}\neq0
\end{equation}
would constitute a particularly strong signature of genuine coherent fission control.

\subsection{Three regimes}

The model identifies three qualitative regimes.

\paragraph{Regime I: no nuclear-state control.}
If
\begin{equation}
W_Q\ll\Gamma_{\mathrm{eq}},
\end{equation}
the externally driven redistribution is negligible and
\begin{equation}
P_f^{(Q)}\simeq P_f^{(0)}.
\end{equation}

\paragraph{Regime II: population control.}
If externally driven redistribution competes with relaxation but coherence is rapidly lost, the density matrix is effectively diagonal and control proceeds through state populations.

\paragraph{Regime III: coherent control.}
If the coherent coupling satisfies approximately
\begin{equation}
\Omega\gtrsim\Gamma_{\mathrm{decoh}}
\end{equation}
over the relevant interval, phase-sensitive interference may become observable.

The second regime is sufficient for the general reactor-control framework and is considerably less demanding than the third.

\subsection{Coupling the microscopic model to a neutron field}

For a two-state population model,
\begin{equation}
\sigma_f^{\mathrm{eff}}(E,t)
=
p_A(t)\sigma_f^{(A)}(E)
+
p_B(t)\sigma_f^{(B)}(E).
\end{equation}

The macroscopic fission rate becomes
\begin{equation}
R_f(t)
=
N_f
\int
\Phi(E,t)
\left[
p_A(t)\sigma_f^{(A)}(E)
+
p_B(t)\sigma_f^{(B)}(E)
\right]
dE.
\label{eq:macro_two_state}
\end{equation}

Thus,
\begin{equation}
u_Q(t)
\rightarrow
p_i(t)
\rightarrow
\sigma_f^{\mathrm{eff}}(t)
\rightarrow
R_f(t).
\end{equation}

\subsection{Effective reactivity contribution}

Suppose external control induces
\begin{equation}
\sigma_f
\rightarrow
\sigma_f+\delta\sigma_f.
\end{equation}

Then
\begin{equation}
\delta\Sigma_f
=
N_f\delta\sigma_f.
\end{equation}

For a sufficiently small perturbation,
\begin{equation}
\delta k_{\mathrm{eff}}
\simeq
\frac{\partial k_{\mathrm{eff}}}{\partial\Sigma_f}
\delta\Sigma_f.
\end{equation}

Since
\begin{equation}
\rho
=
1-\frac{1}{k_{\mathrm{eff}}},
\end{equation}
the associated first-order reactivity change is
\begin{equation}
\delta\rho_Q
\simeq
\frac{\delta k_{\mathrm{eff}}}{k_{\mathrm{eff}}^2}.
\end{equation}

Near criticality,
\begin{equation}
k_{\mathrm{eff}}\simeq1,
\end{equation}
and therefore
\begin{equation}
\delta\rho_Q\simeq\delta k_{\mathrm{eff}}.
\end{equation}

\subsection{Extension of point-reactor kinetics}

Ordinary point kinetics may be written as
\begin{equation}
\frac{dn}{dt}
=
\frac{\rho-\beta}{\Lambda}n
+
\sum_i\lambda_i C_i.
\end{equation}

In the present framework,
\begin{equation}
\rho
=
\rho_{\mathrm{conv}}
+
\rho_Q,
\end{equation}
so that
\begin{equation}
\frac{dn}{dt}
=
\frac{
\rho_{\mathrm{conv}}
+
\rho_Q
-
\beta
}{
\Lambda
}n
+
\sum_i\lambda_iC_i.
\label{eq:extended_point_kinetics}
\end{equation}

Unlike an externally prescribed conventional reactivity insertion, $\rho_Q$ has its own microscopic dynamics:
\begin{align}
\dot{\rho}_N
&=
\mathcal{F}(\rho_N,u_Q,\Phi,\ldots),
\\
\rho_Q
&=
\mathcal{G}(\rho_N),
\\
\dot n
&=
\frac{\rho_{\mathrm{conv}}+\rho_Q-\beta}{\Lambda}n
+
\sum_i\lambda_iC_i.
\end{align}

This system is the minimal mathematical architecture of a multiplying fission system whose fuel possesses an externally controllable internal nuclear-state variable.

\subsection{Feedback}

The coupling may be bidirectional. External actuation modifies the nuclear state,
\begin{equation}
u_Q\rightarrow\rho_N,
\end{equation}
while the neutron field itself can alter nuclear-state populations through absorption and compound-nucleus formation,
\begin{equation}
\Phi\rightarrow\rho_N.
\end{equation}
At the same time,
\begin{equation}
\rho_N
\rightarrow
\sigma_f^{\mathrm{eff}}
\rightarrow
\Phi.
\end{equation}

The resulting structure is therefore a feedback loop rather than a purely feed-forward actuator.

\subsection{Minimal falsifiable prediction}

Consider two otherwise identical experiments on a fissile target exposed to the same neutron distribution,
\begin{equation}
\Phi_1(E)=\Phi_2(E).
\end{equation}

If an external protocol prepares different nuclear-state distributions,
\begin{equation}
\rho_N^{(1)}
\neq
\rho_N^{(2)},
\end{equation}
the hypothesis predicts, under suitable conditions,
\begin{equation}
P_f^{(1)}
\neq
P_f^{(2)}.
\end{equation}

Equivalently,
\begin{equation}
\boxed{
\left.
\frac{\partial P_f}{\partial u_Q}
\right|_{\Phi}
\neq0.
}
\label{eq:falsifiable_prediction}
\end{equation}

Equation~\eqref{eq:falsifiable_prediction} is the simplest operational statement of the proposal. It separates state preparation from direct radiation-induced fission: the incident neutron field is held fixed while an independent actuator modifies the nuclear-state distribution.

\section{Candidate Physical System and Experimental Test Strategy}
\label{sec:candidate_system}

The general framework becomes substantially more concrete when applied to known nuclear isomers. Among presently known systems, the low-lying isomer of $^{235}$U, denoted $^{235m}$U, provides a particularly useful test case because both ingredients required by the proposed mechanism are already partially established: the isomer possesses a fission response different from that of the nuclear ground state, and mechanisms for externally populating the isomer have been experimentally investigated or theoretically proposed.

Importantly, the use of $^{235m}$U in connection with fission control is not itself a new suggestion. Previous work has explicitly noted that the isomer may be exploited to control nuclear fission because its thermal-neutron fission cross section differs substantially from that of ground-state $^{235}$U \cite{DEer1988,LiuWang2022}. The purpose of the present framework is therefore more specific: to treat the externally controllable nuclear-state population as a dynamical control variable, connect it explicitly to an effective fission cross section, and propagate that microscopic modification into macroscopic quantities such as fission rate, neutron multiplication, and reactor reactivity.

\subsection{\texorpdfstring{$^{235m}$U}{235mU} as a candidate system}

The first excited nuclear state of $^{235}$U lies only
\begin{equation}
E_m
=
76.737\pm0.018\ \mathrm{eV}
\end{equation}
above the ground state \cite{Ponce2018}. This exceptionally low excitation energy is unusual for a nuclear transition. The state has a half-life of approximately 26 min and decays predominantly through internal conversion.

For the present proposal, however, the most important property is its different neutron-induced fission response.

D'Eer \textit{et al.} measured the fission cross section of $^{235m}$U relative to that of ground-state $^{235}$U and obtained
\begin{equation}
\frac{\sigma_f^m}{\sigma_f^g}
=
1.61\pm0.44
\end{equation}
for cold neutrons and
\begin{equation}
\frac{\sigma_f^m}{\sigma_f^g}
=
2.47\pm0.45
\end{equation}
for thermal neutrons \cite{DEer1988}.

Thus,
\begin{equation}
\boxed{
\sigma_f^m\neq\sigma_f^g
}
\end{equation}
is experimentally established.

This result demonstrates that the neutron-induced fission response of a fissile isotope is not determined by isotopic identity alone; it can depend strongly on the internal nuclear state.

\subsection{Minimal two-population description}

Let $p_m(t)$ be the fraction of $^{235}$U nuclei occupying the isomeric state and
\begin{equation}
p_g(t)=1-p_m(t)
\end{equation}
the fraction in the ground state.

The effective microscopic fission cross section is
\begin{equation}
\sigma_f^{\mathrm{eff}}(E,t)
=
[1-p_m(t)]\,\sigma_f^g(E)
+
p_m(t)\,\sigma_f^m(E).
\end{equation}

Equivalently,
\begin{equation}
\sigma_f^{\mathrm{eff}}
=
\sigma_f^g
+
p_m
\left(
\sigma_f^m-\sigma_f^g
\right).
\end{equation}

Defining
\begin{equation}
r(E)
=
\frac{\sigma_f^m(E)}{\sigma_f^g(E)},
\end{equation}
one obtains
\begin{equation}
\frac{\Delta\sigma_f}{\sigma_f^g}
=
p_m[r(E)-1].
\end{equation}

For thermal neutrons, using the measured value $r\simeq2.47$,
\begin{equation}
\boxed{
\frac{\Delta\sigma_f}{\sigma_f^g}
\simeq
1.47\,p_m
}.
\label{eq:u235_thermal_effect}
\end{equation}

Equation~\eqref{eq:u235_thermal_effect} exposes the control problem directly. Once the state-dependent fission cross sections are known, the new dynamical quantity is the controllable isomeric fraction $p_m(t)$:
\begin{equation}
u_Q(t)
\longrightarrow
p_m(t)
\longrightarrow
\sigma_f^{\mathrm{eff}}(t).
\end{equation}

The key microscopic feasibility question becomes whether $p_m(t)$ can be externally and reproducibly controlled.

\subsection{External population of the isomer}

Producing a useful controlled population of $^{235m}$U remains the principal unresolved step.

Nuclear excitation by electronic transition (NEET) has been investigated as one possible mechanism. A dedicated search using a laser-generated uranium plasma observed no statistically significant $^{235m}$U production attributable to NEET and established an upper bound under the experimental conditions \cite{Chodash2016}.

Liu and Wang later investigated direct inelastic scattering of low-energy electrons using a distorted-wave treatment \cite{LiuWang2022}. Their calculations predict excitation cross sections larger than earlier plane-wave estimates and also discuss indirect population through higher nuclear levels as a possible route.

These studies do not establish a practical actuator. They instead define the open problem: whether an externally adjustable interaction can generate a sufficiently large, reproducible, and sufficiently rapid modulation of the isomeric fraction.

The present framework does not require a particular excitation technology. What matters is a reproducible mapping
\begin{equation}
u_Q(t)
\rightarrow
p_m(t).
\end{equation}

\subsection{Distinguishing excitation from control}

Producing an isomer is not equivalent to controlling fission through the isomeric population.

The existence of
\begin{equation}
^{235}\mathrm{U}_g
\rightarrow
^{235m}\mathrm{U}
\end{equation}
is not by itself sufficient.

A control mechanism requires
\begin{equation}
p_m=p_m(u_Q)
\end{equation}
and a corresponding modification of neutron-induced fission under an otherwise fixed neutron field:
\begin{equation}
\boxed{
u_Q
\rightarrow
p_m
\rightarrow
\sigma_f^{\mathrm{eff}}
\rightarrow
R_f
}.
\end{equation}

This distinction is what connects nuclear excitation physics to control theory.

\subsection{Minimal falsification experiment}

A decisive test does not require a multiplying reactor.

Consider two otherwise equivalent fissile ensembles exposed to the same neutron spectrum,
\begin{equation}
\Phi_1(E)=\Phi_2(E),
\end{equation}
but with different isomeric fractions,
\begin{equation}
p_m^{(1)}
\neq
p_m^{(2)}.
\end{equation}

The predicted fission-rate difference is
\begin{equation}
\Delta R_f
=
N_f
\int
\Phi(E)
\left[
p_m^{(2)}-p_m^{(1)}
\right]
\left[
\sigma_f^m(E)-\sigma_f^g(E)
\right]
\,dE.
\label{eq:u235_experiment}
\end{equation}

Since the cross-section difference is already known to be nonzero for cold and thermal neutrons, the decisive experimental step is to create
\begin{equation}
\Delta p_m
=
p_m^{(2)}-p_m^{(1)}
\end{equation}
through an externally adjustable actuator.

The hypothesis predicts
\begin{equation}
\boxed{
\left.
\frac{\partial R_f}
{\partial u_Q}
\right|_{\Phi}
\neq0
}
\end{equation}
provided that
\begin{equation}
\frac{\partial p_m}{\partial u_Q}\neq0
\end{equation}
and
\begin{equation}
\sigma_f^m-\sigma_f^g\neq0.
\end{equation}

Observation of a reproducible chain
\begin{equation}
u_Q
\rightarrow
p_m
\rightarrow
R_f
\end{equation}
with the neutron spectrum held fixed would constitute direct evidence of externally actuated nuclear-state control of neutron-induced fission.

\subsection{From microscopic control to macroscopic relevance}

A positive microscopic result would not imply practical reactor control.

Let $f_Q$ denote the fraction of fissile nuclei whose state population can effectively be modified. To first order,
\begin{equation}
\frac{\delta\Sigma_f}{\Sigma_f}
\sim
f_Q
\frac{\delta\sigma_f}{\sigma_f}.
\end{equation}

Macroscopic relevance therefore depends not only on
\begin{equation}
\Delta\sigma_f
=
\sigma_f^m-\sigma_f^g,
\end{equation}
but also on the achievable state population, actuation and relaxation timescales, spatial penetration, and energetic cost.

A minimal feasibility parameter set is
\begin{equation}
\boxed{
\left\{
p_m^{\max},
\frac{dp_m}{dt},
\tau_m,
\Delta\sigma_f(E),
E_Q,
f_Q
\right\}
}.
\end{equation}

Only once these quantities are constrained does it become meaningful to propagate the microscopic effect into $\delta k_{\mathrm{eff}}$ and $\rho_Q$.

\subsection{Scope of the present contribution}

The $^{235m}$U literature requires a careful novelty statement.

Previous studies established that the isomer has a different neutron-induced fission cross section, and the possibility of exploiting this fact for fission control has been explicitly mentioned \cite{DEer1988,LiuWang2022}. The present work therefore does \emph{not} claim the first proposal to use $^{235m}$U to influence fission.

The distinct objective is to formulate nuclear-state population as an explicit dynamical control variable and embed it in the multiscale chain
\begin{equation}
u_Q(t)
\rightarrow
\rho_N(t)
\rightarrow
\sigma_f^{\mathrm{eff}}(E,t)
\rightarrow
R_f(t)
\rightarrow
k_{\mathrm{eff}}(t)
\rightarrow
\rho_Q(t).
\label{eq:multiscale_chain}
\end{equation}

Within this framework, an isomer is not merely a state with unusual reaction properties. Its population becomes a candidate \emph{state variable of the controlled fission system}.

This formulation separates three questions:
\begin{enumerate}
\item \textbf{State dependence:} do different nuclear states have different fission responses?
\item \textbf{Actuation:} can their populations be externally and dynamically controlled?
\item \textbf{System coupling:} can the resulting microscopic modification influence macroscopic neutron multiplication?
\end{enumerate}

For $^{235m}$U, the first question already has a positive experimental answer. The second remains an active nuclear-physics problem. The third is the reactor-physics and control-theory problem formulated here.

\subsection{Beyond \texorpdfstring{$^{235m}$U}{235mU}}

A suitable candidate system should combine two properties:
\begin{equation}
\boxed{
\begin{minipage}{0.82\columnwidth}
\centering
nuclear-state controllability\\
$\times$ fission sensitivity to that state
\end{minipage}
}.
\end{equation}

Possible candidate classes include long-lived spectroscopic isomers, shape isomers, and selected excited configurations of actinides. A state with a large fission effect but negligible controllability is unsuitable; a readily controllable state is equally unhelpful if its fission response is unchanged.

\subsection{Experimental decision tree}

The research program can proceed through increasingly stringent tests:
\begin{enumerate}
\item demonstrate reproducible external control of a fission-distinct nuclear-state population,
\begin{equation}
u_Q\rightarrow\rho_N;
\end{equation}
\item under a fixed neutron spectrum, verify that the controlled state changes neutron-induced fission,
\begin{equation}
\rho_N
\rightarrow
\sigma_f^{\mathrm{eff}}
\rightarrow
R_f;
\end{equation}
\item determine whether the achievable controlled fraction is sufficient to perturb neutron multiplication,
\begin{equation}
R_f
\rightarrow
\delta k_{\mathrm{eff}}.
\end{equation}
\end{enumerate}

This hierarchy separates the fundamental microscopic question from the engineering one. A negative result at the first or second stage would constrain or falsify the proposed mechanism without reactor-scale experiments. A positive result would provide the parameters needed to assess macroscopic relevance.

\section{Outlook: Quantum-State Engineering of Fission Outcomes}
\label{sec:outlook}

The framework developed in the previous sections has focused primarily on the possibility of controlling the probability of fission and, consequently, the macroscopic neutron multiplication of a fissile system. The same conceptual framework suggests a broader question.

A fission event is not characterized by a single final state. For a given fissioning nucleus, the process produces a distribution of fragment masses, charges, excitation energies, neutron multiplicities, and prompt electromagnetic emission. The relevant nuclear response is therefore more general than a single fission cross section.

We introduce the schematic fission-response vector
\begin{equation}
\mathcal{F}
=
\left\{
\sigma_f,\,
Y(A,Z),\,
\nu_n,\,
Y_\gamma,\,
E_{\mathrm{kin}},\,
\ldots
\right\},
\label{eq:fission_response_vector}
\end{equation}
where $Y(A,Z)$ denotes the yield of a fragment with mass number $A$ and atomic number $Z$, $\nu_n$ denotes neutron multiplicity, and the remaining quantities represent other measurable fission observables.

The more general hypothesis is then
\begin{equation}
\boxed{
\mathcal{F}
=
\mathcal{F}(E,\rho_N)
}
\label{eq:fission_response_state}
\end{equation}
rather than only
\begin{equation}
\sigma_f
=
\sigma_f(E,\rho_N).
\end{equation}

If Eq.~\eqref{eq:fission_response_state} can be exploited dynamically, nuclear-state control would extend from \emph{fission-rate engineering} to \emph{fission-outcome engineering}.

\subsection{State dependence of fission-fragment distributions}

Fission-fragment yields are strongly structured rather than statistically featureless. Shell effects, deformation, collective dynamics, excitation energy, angular momentum, and dissipative evolution influence the probability of reaching different scission configurations.

Microscopic and phenomenological calculations reproduce the strong role of shell structure in determining asymmetric fragment distributions \cite{ScampsSimenel2018,Sadhukhan2022}. In actinides, deformed shell effects in the nascent fragments can favor specific asymmetric mass and charge splits \cite{ScampsSimenel2018}.

Experimental results also demonstrate that fission yields are not invariant with respect to the initial conditions of the fissioning system. Measurements of $^{240}$Pu show a systematic evolution of isotopic fission yields with excitation energy, including damping of shell-driven structures and changes in the neutron content of the fragments \cite{Ramos2026}.

These observations establish the general principle
\begin{equation}
\text{initial nuclear conditions}
\longrightarrow
Y(A,Z).
\end{equation}

They do not establish that an external actuator can deliberately select a desired yield distribution. They do, however, show that the fragment distribution contains dynamical information about the state and excitation of the fissioning system.

The natural extension of the present framework is therefore
\begin{equation}
u_Q
\longrightarrow
\rho_N
\longrightarrow
Y(A,Z\,|\,\rho_N).
\label{eq:yield_control_chain}
\end{equation}

\subsection{From state-dependent fissility to state-engineered fissility}

The $^{235m}$U example considered in Sec.~\ref{sec:candidate_system} demonstrates that the neutron-induced fission response of a nucleus can depend strongly on its internal state.

This raises a broader possibility. Instead of classifying a nuclide exclusively through its ground-state fission properties, one may consider a state-dependent criterion,
\begin{equation}
\text{fission response}
=
\text{fission response}(\text{nuclide},\rho_N,E).
\end{equation}

A material that is poorly fissile under a given neutron spectrum in its reference state could, in principle, possess an excited or metastable configuration with a substantially different fission response.

Schematically, one may search for systems satisfying
\begin{equation}
\sigma_f(E;\rho_N^\ast)
\gg
\sigma_f(E;\rho_N^{(0)}),
\label{eq:state_engineered_fissility}
\end{equation}
together with an experimentally accessible pathway
\begin{equation}
\rho_N^{(0)}
\overset{u_Q}{\longrightarrow}
\rho_N^\ast.
\end{equation}

Equation~\eqref{eq:state_engineered_fissility} suggests the concept of \emph{state-engineered fissility}: the possibility that the practical usefulness of a nuclear material as a fission medium may depend not only on isotopic composition but also on the externally prepared nuclear-state distribution.

This possibility remains speculative. In particular, an externally driven excitation is useful as an energy-production control mechanism only if the energetic cost of preparing the relevant state does not simply replace the energy otherwise supplied by the nuclear reaction.

A necessary energetic condition for technological relevance is therefore schematically
\begin{equation}
E_Q
\ll
E_{\mathrm{fission}},
\end{equation}
where $E_Q$ represents the net energy cost required to prepare or maintain the controlled nuclear state per affected fission event.

The external actuator should function primarily as a control or gating mechanism rather than as the dominant energy source of the reaction.

\subsection{Can fission products be steered toward stability?}

A second and potentially more consequential extension concerns the radioactive character of the fission products.

For conventional actinide fission, the primary fragments are generally neutron rich. They typically undergo prompt neutron emission followed by sequences of beta decays before reaching stable or long-lived nuclei. Therefore, requiring a fission event to produce two immediately stable fragments is a much stronger condition than merely modifying the fragment-yield distribution.

Conservation of nucleon number, charge, energy, and angular momentum strongly constrains the accessible final states. Nuclear-state control cannot remove these constraints.

The scientifically meaningful question is therefore not initially whether every fission can produce stable fragments, but rather whether nuclear-state preparation can bias the fission-yield distribution toward products closer to stability or away from long-lived radionuclides.

This problem can be expressed quantitatively.

Let
\begin{equation}
Y(A,Z\,|\,\rho_N)
\end{equation}
denote the independent yield distribution produced from an initial nuclear state $\rho_N$.

Define a weight
\begin{equation}
w(A,Z)
\end{equation}
that represents an undesired property of each fragment, such as long half-life, long-term radiotoxicity, delayed-neutron contribution, decay heat, or another application-specific quantity.

An objective functional may then be introduced:
\begin{equation}
\mathcal{J}[\rho_N]
=
\sum_{A,Z}
w(A,Z)
Y(A,Z\,|\,\rho_N).
\label{eq:waste_functional}
\end{equation}

The fission-outcome control problem becomes
\begin{equation}
\boxed{
\rho_N^{\mathrm{opt}}
=
\operatorname*{arg\,min}_{\rho_N}
\mathcal{J}[\rho_N]
}
\label{eq:optimal_state}
\end{equation}
subject to physical accessibility of $\rho_N$, conservation laws, fission probability, and the energetic cost of the actuator.

Equation~\eqref{eq:optimal_state} converts the qualitative idea of ``cleaner fission'' into a falsifiable optimization problem.

For example, $w(A,Z)$ may be chosen to penalize selected long-lived fission products. Alternatively, one could minimize a time-dependent activity or radiotoxicity measure,
\begin{equation}
\mathcal{J}(t;\rho_N)
=
\sum_{A,Z}
w(A,Z,t)
Y(A,Z\,|\,\rho_N).
\end{equation}

The central microscopic quantity required to assess this possibility is then
\begin{equation}
\frac{\partial Y(A,Z\,|\,\rho_N)}
{\partial u_Q},
\end{equation}
or, more fundamentally,
\begin{equation}
\frac{\delta Y(A,Z)}
{\delta\rho_N}.
\end{equation}

\subsection{Two complementary control objectives}

The nuclear-state engineering framework therefore naturally separates into two distinct control objectives.

The first is \emph{reactivity control}:
\begin{equation}
u_Q
\rightarrow
\rho_N
\rightarrow
\sigma_f^{\mathrm{eff}}
\rightarrow
R_f
\rightarrow
k_{\mathrm{eff}}.
\label{eq:reactivity_objective}
\end{equation}

The second is \emph{fission-outcome control}:
\begin{equation}
u_Q
\rightarrow
\rho_N
\rightarrow
Y(A,Z),
\ \nu_n,
\ Y_\gamma,
\ldots
\label{eq:outcome_objective}
\end{equation}

The two objectives need not require the same nuclear states or the same actuator.

A state that produces a large change in $\sigma_f$ may leave the fragment distribution almost unchanged. Conversely, a state that produces only a modest change in total fission probability could still significantly alter the relative probability of competing fission pathways.

This distinction implies that nuclear-state engineering could remain useful even if the induced reactivity perturbation is too small to replace conventional reactor-control systems. A small but reproducible shift in selected fragment yields could have scientific or technological value independently of reactor-power control.

\subsection{A broader research program}

The framework suggests a hierarchy of increasingly ambitious questions:

\begin{enumerate}
\item Can an external actuator reproducibly control the population of a nuclear state with a different neutron-induced fission response?
\item Can such control produce a measurable modification of the total fission probability under a fixed neutron field?
\item Can nuclear-state preparation alter fragment yields or other fission observables?
\item Can an initial state be identified that suppresses selected undesirable fission products?
\item Can useful state-dependent fission properties be found in nuclides not normally regarded as practical fissile fuels under the same neutron conditions?
\end{enumerate}

The first two questions define the core proposal developed in this paper. The remaining questions constitute a broader research program that may be termed
\begin{equation}
\boxed{
\begin{minipage}{0.82\columnwidth}
\centering
\textit{quantum-state engineering of nuclear fission}.
\end{minipage}
}
\end{equation}

The term is used here in a precise sense: the deliberate preparation of nuclear-state distributions with the objective of modifying one or more measurable properties of a subsequent fission process.

No assumption is made that arbitrary fission products can be selected, that stable fragments can necessarily be produced directly, or that such control can be achieved at technologically relevant scales. Those are quantitative questions to be determined by nuclear structure, reaction dynamics, and experiment.

The conceptual extension is that the initial nuclear state is treated as a potentially engineerable input to the fission process rather than solely as an inaccessible microscopic variable encoded in averaged reaction data.

\section{Conclusions}
\label{sec:conclusions}

Nuclear reactor physics conventionally incorporates microscopic quantum phenomena through reaction data while treating the internal nuclear state of the fuel as an effectively uncontrolled variable. This work has explored a generalization in which the nuclear state is promoted to an explicit dynamical variable,
\begin{equation}
\rho_N(t),
\end{equation}
and the effective fission response becomes
\begin{equation}
\sigma_f
=
\sigma_f(E,\rho_N).
\end{equation}

If an external actuator can modify the nuclear-state distribution,
\begin{equation}
u_Q
\rightarrow
\rho_N,
\end{equation}
then a complementary control pathway can be defined:
\begin{equation}
u_Q
\rightarrow
\rho_N
\rightarrow
\sigma_f^{\mathrm{eff}}
\rightarrow
R_f
\rightarrow
k_{\mathrm{eff}}.
\end{equation}

The individual physical ingredients of this chain are not entirely hypothetical. Nuclear quantum-optical experiments have demonstrated externally controlled nuclear excitation dynamics in selected systems \cite{Heeg2021,Yamaguchi2024}. More directly relevant to fission, measurements of the low-lying $^{235}$U isomer demonstrate that a change in nuclear state can produce a substantial change in neutron-induced fission cross section \cite{DEer1988}.

Thus, for a relevant fissile isotope, the relation
\begin{equation}
\rho_N
\rightarrow
\sigma_f
\end{equation}
is experimentally established.

What remains unresolved is the controllability of the first step,
\begin{equation}
u_Q
\rightarrow
\rho_N,
\end{equation}
at an efficiency, timescale, and addressed fraction large enough to matter.

Several mechanisms for populating $^{235m}$U have been investigated, including nucleus-electron coupling and inelastic electron scattering \cite{Chodash2016,LiuWang2022}, but no presently established method provides efficient and tunable control of a macroscopic isomeric population suitable for reactor applications.

The central physical question is therefore not simply whether fission probability depends on nuclear state. For $^{235}$U, it demonstrably does. The relevant questions are instead
\begin{equation}
\boxed{
\begin{minipage}{0.86\columnwidth}
\centering
Can the population of fission-distinct nuclear states be externally\\
and dynamically controlled?
\end{minipage}
}
\end{equation}
and, subsequently,
\begin{equation}
\boxed{
\begin{minipage}{0.86\columnwidth}
\centering
Can such microscopic control produce a non-negligible\\
macroscopic reactivity effect?
\end{minipage}
}
\end{equation}

The minimal two-state model shows how these questions can be expressed quantitatively. For the $^{235}$U ground-state/isomer pair,
\begin{equation}
\sigma_f^{\mathrm{eff}}
=
(1-p_m)\sigma_f^g
+
p_m\sigma_f^m,
\end{equation}
which provides a direct connection between a nuclear-state population and reactor quantities.

At the system level, the framework can be represented by an additional perturbative reactivity contribution,
\begin{equation}
\rho
=
\rho_{\mathrm{conv}}
+
\rho_Q,
\end{equation}
with
\begin{equation}
\rho_Q
=
\mathcal{G}[\rho_N].
\end{equation}

This formulation should be understood as a theoretical control framework, not as an immediately realizable reactor technology. State-preparation efficiency, the controllable fraction of nuclei, penetration and spatial uniformity of the actuator, relaxation and decoherence, competing reaction channels, energy requirements, and coupling to neutron kinetics remain major open problems.

The literature also constrains the novelty claim. State-dependent fission in $^{235m}$U is known, and prior work has already mentioned its possible use for fission control \cite{LiuWang2022}. The contribution proposed here is therefore narrower: the explicit formulation of nuclear-state population as a dynamical control variable coupled through a multiscale chain to fission rate, neutron multiplication, and reactor reactivity.

In this formulation, conventional and nuclear-state control operate on complementary parts of the fission process:
\begin{equation}
\boxed{
\begin{minipage}{0.82\columnwidth}
\centering
conventional control $\rightarrow$ control of the neutron environment
\end{minipage}
}
\end{equation}
whereas
\begin{equation}
\boxed{
\begin{minipage}{0.82\columnwidth}
\centering
nuclear-state control $\rightarrow$ control of the nuclear response of the fuel.
\end{minipage}
}
\end{equation}

The immediate next step is not the design of a nuclear-state-controlled reactor. It is the identification, calculation, and experimental study of systems in which an externally adjustable actuator produces a measurable change in a fission-distinct nuclear-state population and hence a measurable change in neutron-induced fission under an otherwise fixed neutron field.

The operational criterion is
\begin{equation}
\boxed{
\left.
\frac{\partial R_f}
{\partial u_Q}
\right|_{\Phi}
\neq0.
}
\end{equation}

A verified nonzero response would supply the microscopic parameters required to determine whether nuclear-state control can remain a laboratory-scale nuclear effect or become a meaningful new degree of freedom in the dynamics of multiplying fission systems.


\begin{acknowledgments}
The authors thank F. Sisini for a private communication that
motivated the initial formulation of the problem. The final form
of the manuscript emerged from an independent exploration of the
consequences of that idea.
\end{acknowledgments}

\bibliography{references}

\section*{Fictional Work Notice}

This document is a fictional scientific article created exclusively
for narrative and artistic purposes.

The authors named in this manuscript are fictional. Any resemblance
to actual persons, living or deceased, is entirely coincidental and
unintentional.

Real institutions, facilities, scientific publications, and technical
terminology may be referenced in order to provide a realistic setting.
Their appearance in this fictional work does not imply participation,
endorsement, authorship, or approval by those institutions or by their
personnel.

The scientific scenario described in this manuscript forms part of a
work of fiction and should not be interpreted as reporting actual
experiments, results, technologies, or research programs.

\end{document}
